\selectlanguage{english}

\subsection*{Introduction}
This bachelor thesis examines whether a practical Linux \gls{cli} assistant can
be built on small local \glspl{llm} without relying on cloud services. The
motivation is to combine the usability of natural-language assistance with the
privacy, offline availability, and low-latency benefits of local execution. The
work targets resource-constrained hardware and explicitly evaluates which
architectural approaches remain reliable under such constraints.

The thesis sets three primary goals: (1) to analyze the landscape of existing
terminal assistants and identify their limitations for local use, (2) to design
and implement a prototype system that operates entirely on local models, and (3)
to evaluate the feasibility of key components using objective metrics. The
resulting prototype, shAI, is open source and serves as a reference for future
research on constrained \gls{llm} systems.

The document is organized to mirror this workflow. It begins with theoretical
background, proceeds to a systematic review of existing solutions, then presents
the design and implementation of shAI, and finally evaluates each component with
measured results. The resume below follows the same per-section structure and
condenses each chapter into its core contributions and findings.

Two research questions guide the work: which \gls{llm}-based tasks remain feasible
on small local models, and how should system architecture be adapted to avoid
hallucination and latency spikes. The thesis treats these questions as empirical
problems and emphasizes repeatable evaluation. This focus on measurement leads to
the creation of datasets and test configurations that allow direct comparison
between models, retrieval strategies, and latency trade-offs.

\subsection*{Theoretical Background}
The theoretical section begins with terminal interaction concepts. It defines
terminal emulators, \glspl{cli}, and \glspl{tui}, and describes how shells and
command-line utilities form a text-based interface for user interaction. These
concepts establish the environment in which the assistant must operate and the
constraints imposed by standard input/output streams.

The thesis then summarizes \gls{llm} architecture, including tokenization,
embedding layers, and transformer blocks. It distinguishes dense models from
\gls{moe} variants and highlights their trade-offs in memory usage, computational
efficiency, and output consistency. Key limitations of local models are
emphasized: smaller parameter counts reduce reasoning ability, and long prompts
increase latency and hallucination risk.

Additional considerations include tokenizer limitations for shell syntax and
flags, which can confuse small models even on short prompts. The section also
discusses how context window size affects prompt engineering and why iterative
processing can be necessary when documents exceed practical limits. This theory
provides a basis for later design decisions around passage selection, query
reformulation, and output length control.

Embedding models and document representation are introduced as the foundation
for retrieval. The text explains dense vs.\ sparse retrieval, semantic vector
representations, and vector databases, with \gls{faiss} selected for its
lightweight file-oriented indexing. The \gls{rag} pipeline is described as a
two-stage process of retrieval followed by summarization, and retrieval
granularity is compared (document-level vs.\ passage-level). Finally, the
evaluation metrics used in the thesis are defined, including Precision@k,
Recall@k, Mean Reciprocal Rank (MRR), and \gls{ndcg}, along with query type
taxonomy for \gls{rag} evaluation.

Additional background covers practical \gls{rag} concerns: prompt size overhead,
iterative prompt processing, and document relevance filtering. The thesis adopts
the query-type framework (levels 1--4) to distinguish explicit fact retrieval
from tasks that require hidden rationales or advanced reasoning. The focus is on
level-1 and level-2 queries, which align with man page assistance and can be
supported by retrieval and lightweight summarization.

The evaluation subsection elaborates on how retrieval metrics map to user
experience. Precision@k and Recall@k capture the breadth of relevant results,
while MRR emphasizes early ranking for quick responses. \gls{ndcg} is introduced
as a graded relevance metric even though it is not used in the final tests,
helping position the chosen metrics among standard information-retrieval
benchmarks. This theoretical framing sets the stage for the empirical evaluation
chapters.

The section also describes Linux man pages as structured documents, including
their conventional sections (NAME, SYNOPSIS, DESCRIPTION, OPTIONS, and others)
and the organization of pages into manual groups. This context informs the
document processing pipeline and defines what constitutes a document or passage
in later experiments.

Document representation and storage are discussed in more detail, including dense
and sparse retrieval approaches and the role of vector databases. The thesis
explains why a lightweight, file-based vector index is suitable for a prototype
and outlines common embedding aggregation methods (pooling, attention-based
summaries, and special-token representations). These foundations justify the use
of embedding \glspl{llm} for semantic search in the later implementation.

The thesis also summarizes the model families used in experiments and their
relevant characteristics. Qwen2.5 and Qwen3 models provide small parameter
variants with long context windows, Granite offers dense and \gls{moe} variants
optimized for efficiency, and Gemma models serve as a lightweight baseline. The
embedding models mirror this diversity, ranging from small Granite embeddings to
Qwen3 embeddings that support larger dimensional outputs. These choices reflect
the requirement to test performance on realistic local hardware.

The retrieval evaluation metrics are also linked to practical user behavior. For
example, Precision@1 captures whether the first result is immediately useful,
while Precision@5 approximates the ability to find a relevant answer in a small
set of results. MRR rewards early correct hits and is especially relevant for
interactive assistants that should not require extensive browsing. These
connections justify the choice of metrics for the experimental section.

\subsection*{Analysis of Existing Solutions and Technologies}
The thesis surveys existing \gls{cli} and terminal assistants, including shell
companions, terminal emulators with AI features, and agentic \gls{cli} tools.
Solutions such as Butterfish, llm-zsh-plugin, Warp Terminal, gptme, gemini-cli,
AiTerm, and others are compared using criteria that matter for local use:
open-source availability, local \gls{llm} support, context awareness, agentic
architecture, tool integration, and implementation language.

The analysis finds that most mature tools depend on remote APIs or proprietary
services, while those that support local models struggle with long prompts,
hallucinations, and tool reliability. Local tests with small models show that
agentic workflows are fragile and often fail due to incorrect tool invocation or
multi-step reasoning breakdowns. These findings motivate a scope shift away from
autonomous shell execution and toward a constrained information-retrieval task
that is more compatible with local model capabilities.

Specific observations include weak local performance of agentic tools when
running Gemma and Qwen models, where long prompts lead to hallucinated steps and
incorrect tool usage. In contrast, tools that focus on retrieval or short-form
generation provide more stable results. The analysis also notes that terminal
emulators with AI features, such as Warp, provide a rich user experience but are
closed source and not aligned with local privacy constraints, making them
unsuitable as a foundation for this work.

The survey also yields a structured comparison table across criteria such as
open-source availability, local model support, context awareness, and tool
integration. This comparison highlights that remote services dominate the space,
and local-capable tools tend to trade reliability for broad functionality. The
thesis documents a three-step scope evolution: an initial vision of a broad
shell wrapper, a local-first design with many features, and a final, focused
\gls{rag} architecture that was empirically testable within the thesis timeline.

Several features were explicitly deprioritized after this analysis. Seamless
shell integration would require shell-specific hooks (Bash, Zsh, Fish), and even
basic command execution complicates testing and validation. Dangerous command
detection was also dropped because it depends on user-specific context and would
require a dedicated labeled dataset. These constraints reinforced the decision to
prioritize retrieval and summarization over open-ended automation.

\subsection*{Design and Implementation}
The design requirements emphasize offline operation, minimal external
dependencies, and a predictable \gls{cli} user experience. Linux man pages are
chosen as the primary knowledge source because they are structured, locally
available, and relevant for command-line assistance. To reduce noise and improve
retrieval quality, man pages are converted to plain text before indexing.

The shAI system is implemented as a modular monolithic application with three
components: (1) a command-vs.-query classifier, (2) a document retriever that
builds and searches embedding indices, and (3) a summarizer that produces concise
responses from retrieved content. The classifier explores two approaches: an
\gls{llm}-only classifier and a hybrid method that combines syntax checking with
\gls{llm} validation. The retriever builds vector indices using local embedding
models via Ollama and \gls{faiss}, and it exposes three retrieval variants:
document-level retrieval, document retrieval with passage-level re-ranking, and
passage-level retrieval based on fixed-size line windows.

The retrieval pipeline includes explicit configuration parameters for top-k
results, passage size, and index storage location, enabling controlled
experiments. Indexing is performed on preprocessed man pages stored as text files
and cached in a local database directory. Search returns both metadata and
content segments, which are then merged or filtered before being passed to the
summarizer.

Beyond the core components, the implementation includes user-facing workflows
for system initialization, index building, query execution, and configuration
changes. The documented use cases include starting the Ollama runtime, building
the embedding index, running search queries, and generating answers. These steps
are reflected in sequence diagrams that describe how the system transitions from
raw query input to final response generation.

The classifier component is designed as a proof of concept. It can be invoked on
the current user input buffer and return a binary decision, but integrating this
into a general-purpose shell requires different hooks for Bash, Zsh, and Fish.
The thesis documents a Bash-oriented approach with readline hooks and explains
why full multi-shell support is out of scope. This discussion complements the
evaluation results, which already show insufficient accuracy for reliable
deployment.

The retriever component is the most mature part of the system. It handles
document ingestion, embedding generation, and vector index storage, and it
supports both document-level and passage-level search. For passage-level search,
documents are split into fixed line windows to control context size. The
re-ranking variant performs a second-stage search within top-k documents, which
helps isolate relevant passages but increases latency.

The summarizer component provides the final user-facing response. It receives a
query and a set of retrieved passages, formats a prompt with explicit structure,
and generates a concise answer focused on the user question. This keeps the
generation task narrow and reduces hallucination compared to open-ended
conversational assistants.

The summarizer orchestrates the \gls{rag} flow: it sends the user query to the
retriever, selects relevant passages, and constructs a compact prompt for the
generative model. The implementation prioritizes small prompts and explicit
structure to reduce hallucination risk. Deployment uses a Dockerized Ollama
runtime with OpenAI-compatible APIs, enabling model swaps and configuration via
simple config files. Service-check logic ensures the local \gls{llm} endpoint is
available before processing requests.

Design decisions are supported by architecture and use-case diagrams: the
component model highlights the separation between retrieval, classification, and
summarization, while the use-case diagram outlines system initialization, index
building, querying, and configuration. The retriever includes two performance
optimizations: document relevance filtering based on short descriptions and
iterative prompt processing that evaluates documents in smaller chunks rather
than injecting entire manuals into a single prompt.

The technology stack centers on Python with the OpenAI-compatible client library
for model access, \gls{faiss} for similarity search, and standard system
libraries for filesystem and subprocess management. Ollama provides local model
runtime with configurable context length and GPU support. The test suite is
config-driven, enabling reproducible runs across model variants and retrieval
strategies.

Implementation also includes a CLI entry point that exposes configuration loading
and test execution. Configuration files specify which model to use for embedding
and generation, and the system is designed to reuse the same interfaces for both
local and remote backends if needed. This modular structure allows the retrieval
and summarization components to be tested independently.

\subsection*{Evaluation and Key Findings}
Evaluation is component-focused and implemented in the test-ctf suite. The
classifier is tested on \texttt{classifier-100}, a balanced dataset of 100 inputs
split equally between natural language questions and shell commands. Retrieval is
tested on \texttt{sample-10} (14 documents, 250 queries) and \texttt{sample-100}
(101 documents, 300 queries), using Precision@k and MRR to evaluate accuracy and
per-query latency for performance.

The classifier results show a clear trade-off between accuracy and latency. The
best model achieved 88\% accuracy but required 2.9s per query, while smaller
models stayed near chance-level (50--65\%) with lower latency. This makes
real-time command classification impractical for local \glspl{llm} in the tested
range.

Per-model results reinforce this conclusion. Granite 1b reached 53\% accuracy,
Granite 3b achieved 60\%, Qwen3 0.6b reached 79\%, Qwen3 1.7b reached 88\%, Qwen2.5
0.5b achieved 50\%, Qwen2.5 1.5b achieved 78\%, and Gemma3 1b achieved 65\%. The
best accuracy coincided with the highest latency, confirming that the
classification approach is not competitive for interactive use.

Models used in classifier testing include Granite (1b and 3b), Qwen2.5
(0.5b/1.5b), Qwen3 (0.6b/1.7b), and Gemma3 1b, with accuracies spanning 50--88\%
and average latencies ranging from 0.10s to 2.9s per query. These results
consistently favor retrieval-based workflows over classification-heavy designs.

Document-level retrieval achieved strong performance. On \texttt{sample-10},
Top-1 accuracy ranged from 83.6\% to 86.8\% with Top-5 up to 100\% and MRR up to
0.916. On \texttt{sample-100}, Top-1 ranged from 62.7\% to 72.7\% with Top-5 up to
92.0 and MRR up to 0.806. Granite embeddings produced the lowest per-query
latencies (about 0.04--0.10s), while qwen3-embedding increased index build time
dramatically on larger corpora (128--134s vs.\ 1.6--2.3s for Granite models).

Document retrieval with passage-level re-ranking increased latency (0.6--3.8s per
query) and reduced Top-1 accuracy (39.2--55.2\% on \texttt{sample-10} and
42.7--45.7\% on \texttt{sample-100}) without improving Top-5 results. Passage-level
retrieval matched document-level precision while keeping low latency for Granite
embeddings (0.03--0.10s). These results indicate that passage-level retrieval can
deliver similar accuracy without the overhead of re-ranking.

Across embedding models, Granite variants offered the best latency and stable
accuracy, while qwen3-embedding provided strong accuracy but with high index
build time on larger corpora. MRR values ranged from 0.861 to 0.916 on
\texttt{sample-10} and from 0.724 to 0.806 on \texttt{sample-100}. These numbers
confirm that top-ranked results are usually relevant but that broader corpora
reduce precision, reinforcing the need for focused retrieval strategies.

The evaluation also highlights the effect of retrieval strategy on index build
time. Granite embeddings built indices in roughly 1.0--2.3s for the tested
datasets, while qwen3-embedding required 10.3s on \texttt{sample-10} and
128--134s on \texttt{sample-100}. These costs are significant for iterative
development and suggest that smaller embeddings are preferable when frequent
re-indexing is expected.

All retrieval tests used top\_k=5, with re-ranking configured to top\_k\_extended=10
and passage\_size=20 lines. These parameters reflect a balance between accuracy
and latency and were kept constant across datasets to enable fair comparisons.

Summarization was evaluated qualitatively rather than with a dedicated benchmark.
The summarizer produced concise and useful answers when retrieval was accurate,
but its output quality degraded when retrieved passages were noisy or incomplete.

The evaluation artifacts are stored as structured JSON outputs, enabling
inspection of per-query ranks, latencies, and match statistics. This approach
supports reproducibility and makes it possible to compare retrieval strategies
under the same configuration. However, the thesis does not yet include an
end-to-end benchmark that combines retrieval and summarization into a single
quality metric, leaving full-system evaluation as future work.

\subsection*{Conclusion and Future Work}
The thesis concludes that local \gls{llm} assistants are viable when constrained
to retrieval and summarization tasks, while classification-heavy or agentic
architectures remain unreliable on small models. The evidence shows that
classification accuracy peaks at 88\% with high latency, whereas retrieval
achieves strong Top-5 accuracy (92--100\%) with low latency when using smaller
embedding models.

The work contributes an open-source prototype, a reproducible evaluation suite,
and empirical findings that guide future development. Future work should focus
on model optimization (quantization and distillation), richer yet lightweight
user interfaces, and a dedicated summarization-quality benchmark. Additional
constrained tasks, such as structured error explanation or template-based command
generation, are promising directions that align with the strengths of small local
\glspl{llm}.

The thesis also identifies clear limitations: the absence of end-to-end
summarization benchmarks, the lack of automated evaluation for agentic behaviors,
and the dependence of answer quality on retrieval precision. Addressing these
limitations will require new datasets, more rigorous evaluation protocols, and
careful balancing between latency and accuracy for local deployment.
